\partdivider{A}{The Nucleus and Its Decay}{A nucleus weighs less than its parts, and that missing mass, the binding energy, decides which nuclei are stable, which decay, and how much energy each decay releases.}{\chref{ch:nucleus} The Nucleus\\[1.2mm]
\chref{ch:decay} Radioactive Decay\\[1.2mm]
\chref{ch:decaylaw} The Law of Radioactive Decay}

\ncertchapter{1}{The Nucleus}\chlabel{ch:nucleus}

Every atom has at its centre an object tens of thousands of times smaller than the atom itself, which
carries all of its positive charge and more than \num{99.9} percent of its mass. Chemistry never sees this
object: chemical reactions rearrange electrons and leave the nucleus untouched. Yet the nucleus holds
the deepest energy store in ordinary matter. A gram of radium keeps itself warmer than its
surroundings for centuries; a kilogram of uranium-235, completely fissioned, gives as much heat as
burning about three thousand tonnes of coal; and the Sun has been
shining on nuclear energy for four and a half billion years. To understand any of this, one needs to know
what a nucleus is made of, how big it is, and how tightly its parts are held.

This chapter answers those three questions. It begins with the experiment that found the nucleus,
which uses nothing beyond the energy conservation of \chref{ch:energy} and the angular momentum of
\chref{ch:rotation}, and it extracts the size of the nucleus from that experiment and its successors.
It then meets the one genuinely new principle of this volume, that mass is a form of energy, and
uses it to define the \emph{binding energy} of a nucleus: the energy that would be needed to take it
apart. A single graph, the binding energy per nucleon plotted against the number of nucleons, then
explains in one stroke why heavy nuclei release energy when they split and light ones when they fuse,
and it will be the reference point for every later chapter. The chapter closes with the force that
does the binding, and with the energy needed to remove a single nucleon.

The limits are worth stating at the start. The nucleus will be treated as a uniform sphere of
nucleons with a sharp edge, and its binding energy will be modelled by a liquid drop. Both pictures
are averages: real nuclei have a soft surface about two femtometres thick, and the lightest nuclei
are better described by the shell structure that the liquid drop ignores. Each place where the model
stops working is flagged as it arises.

\section{Finding the Nucleus and Naming It}\label{sec:nuc-naming}

Before 1909 the accepted picture of the atom was J.\,J.~Thomson's ``plum pudding'': positive charge
spread thinly through the whole volume of the atom, about \SI{e-10}{m} across, with electrons embedded
in it like plums. Between 1909 and 1913 Hans Geiger and Ernest Marsden, working in Ernest Rutherford's
laboratory, fired $\alpha$ particles (fast helium ions, of charge $+2e$ and mass about four times the
hydrogen atom's) at thin metal foils. Almost all of them passed through with tiny deflections, but
about one in eight thousand bounced back through more than $90^\circ$. A thin smear of positive charge
cannot do that. The electric field inside a uniformly charged sphere of atomic size is far too weak to
turn back an $\alpha$ particle moving at five percent of the speed of light, and electrons are thousands
of times too light to do it (a bowling ball is not reversed by ping-pong balls). Rutherford concluded
in 1911 that the positive charge and nearly all the mass are concentrated in a tiny central
\emph{nucleus}; only near such a concentrated charge is the electric force strong enough to reverse
the particle.

Later experiments identified what the nucleus is made of: the \emph{proton} (Rutherford, 1919), the
nucleus of the hydrogen atom, with charge $+e$; and the electrically neutral \emph{neutron}
(Chadwick, 1932), slightly heavier than the proton. Together they are called \emph{nucleons}. A
nucleus is described by how many of each it contains.

\begin{defbox}{Nuclide notation}
A \emph{nuclide} is a nucleus with a definite number of protons and neutrons. The number of protons
$Z$ is the \emph{atomic number}; it fixes the nuclear charge $Ze$ and hence the chemical element. The
total number of nucleons $A$ is the \emph{mass number}, so the number of neutrons is $A-Z$. The nuclide
is written
\[
\nuc{A}{Z}{X},
\]
where X is the chemical symbol. Since the symbol already fixes $Z$, the shorter form $\iso{A}{X}$ (read
``X-$A$'', as in uranium-238) is just as common. Both $A$ and $Z$ are pure numbers.
\end{defbox}

Figure~\ref{fig:nuc-notation} reads one symbol in full. In this book the neutron number is always
written as $A-Z$; the letter $N$ is kept for the number of undecayed nuclei in a sample, which is the
central quantity of \chref{ch:decaylaw}.

\begin{figure}[htbp]
\centering
\tikzfig{nuc_notation}
\figcaption{fig:nuc-notation}{Reading the symbol $\nuc{238}{92}{U}$. The upper number counts all
nucleons, the lower counts protons, and their difference is the number of neutrons. The lower number
is redundant once the symbol is known, which is why it is often dropped.}
\end{figure}

Three families of related nuclides recur so often that they have names (Table~\ref{tab:nuc-families}).
\begin{tablebox}{Isotopes, isobars and isotones}
\centering
\begin{tabular}{@{}llll@{}}\toprule
\textbf{Name} & \textbf{Same} & \textbf{Different} & \textbf{Example}\\\midrule
isotopes & $Z$ & $A-Z$ and $A$ & $\iso{12}{C}$, $\iso{13}{C}$, $\iso{14}{C}$\\
isobars & $A$ & $Z$ and $A-Z$ & $\iso{14}{C}$, $\iso{14}{N}$\\
isotones & $A-Z$ & $Z$ and $A$ & $\iso{13}{C}$, $\iso{14}{N}$ (seven neutrons each)\\\bottomrule
\end{tabular}\label{tab:nuc-families}
\end{tablebox}
Isotopes are chemically identical, because chemistry depends only on the electron cloud, and the
electron cloud depends only on $Z$. This is why a radioactive isotope can stand in for a stable one
inside a living molecule and act as a tracer, one of the applications of \chref{ch:nucapps}. Isobars
matter because, as \chref{ch:decay} shows, $\beta$ decay changes $Z$ while leaving $A$ fixed, so it
moves a nucleus along a row of isobars.

\section{How Big Is a Nucleus?}\label{sec:nuc-size}

Rutherford's experiment does more than show that the nucleus exists: it puts an upper limit on its
size. The reasoning needs only the conservation laws of Volume~I.

\subsection{Head-on: the distance of closest approach}
An $\alpha$ particle aimed straight at a nucleus slows down as it climbs the electric potential-energy
hill, stops momentarily, and is pushed back. At the turning point all of its kinetic energy has become
electric potential energy. The electric potential energy of two charges $q_1$ and $q_2$ a distance $r$
apart is $q_1q_2/(4\pi\varepsilon_0 r)$ (from Class~12 electrostatics; $\varepsilon_0$ is the
permittivity of free space), and for nuclear work the combination
\[
\frac{e^2}{4\pi\varepsilon_0}=\SI{1.44}{\MeV.\fm}
\]
is the convenient form: two elementary charges $r$ femtometres apart ($\SI{1}{\fm}=\SI{e-15}{m}$) have
potential energy $1.44/r$ in \si{\MeV}. The electron-volt and its multiples are defined in the
units table of the volume opening; $\SI{1}{\MeV}=\SI{1.602e-13}{J}$.

\begin{derivbox}{Distance of closest approach}
An $\alpha$ particle (charge $2e$) of kinetic energy $K$, far from a nucleus of charge $Ze$, moves
straight towards it. Treat the nucleus as fixed (its recoil is discussed below). The only force is
the electric repulsion, which is conservative, so $K+U$ is constant (\chref{ch:energy}). Far away
$U=0$; at the turning point, a distance $d$ from the centre, the particle is momentarily at rest. Hence
\begin{equation}\label{eq:nuc-closest}
K=\frac{1}{4\pi\varepsilon_0}\frac{(2e)(Ze)}{d}
\qquad\Longrightarrow\qquad
d=\frac{2Ze^2}{4\pi\varepsilon_0 K}=\frac{2Z\,(\SI{1.44}{\MeV.\fm})}{K}.
\end{equation}
For a \SI{5.3}{\MeV} $\alpha$ particle (from polonium-210) and gold ($Z=79$),
$d=2(79)(1.44)/5.3=\SI{43}{\fm}$.
\end{derivbox}

The angular distribution of Rutherford's scattered particles matched, in every detail, the prediction for
a point charge. So the $\alpha$ particles never touched the gold nucleus, and its radius must be
\emph{less} than about \SI{43}{\fm}, some three thousand times smaller than the atom. Treating the gold
nucleus as fixed is a good approximation: it is 49 times heavier than the $\alpha$ particle, and
\chref{ch:collisions} shows that a heavy target takes only a small fraction of the energy.

\subsection{Off-centre: the impact parameter}
Most $\alpha$ particles are not aimed exactly at a nucleus. The sideways offset of the incoming line
from the centre of the nucleus is called the \emph{impact parameter} $b$. Figure~\ref{fig:nuc-rutherford}
shows computed paths for several values of $b$: the closer the aim, the larger the deflection, and
only a nearly head-on particle comes back. How close does an off-centre particle get?

\begin{figure}[htbp]
\centering
\tikzfig{nuc_rutherford}
\figcaption{fig:nuc-rutherford}{Computed paths of \SI{7.7}{\MeV} $\alpha$ particles passing a gold
nucleus at impact parameters $b$ from \SIrange{8}{100}{\fm}. Each path is a hyperbola with the nucleus at
its outer focus. Notice how small $b$ must be for a large deflection: only particles aimed within about
\SI{15}{\fm} of the centre are turned back by more than $90^\circ$. Even the head-on particle stops at
\SI{30}{\fm}, four times the radius of the gold nucleus (the black disc), so Rutherford's particles
felt only the electric force.}
\end{figure}

\begin{derivbox}{Closest approach at impact parameter \texorpdfstring{$b$}{b}}
The electric force on the $\alpha$ particle always points along the line from the nucleus, so it
exerts no torque about the nucleus and the particle's angular momentum about the nucleus is conserved
(\chref{ch:rotation}; it is the same argument that gives Kepler's second law in \chref{ch:kepler}).
Far away, the particle of mass $m$ moves with speed $v_0$ along a line a distance $b$ from the nucleus,
so its angular momentum is $mv_0b$. At closest approach, a distance $r_{\min}$ from the nucleus, the
velocity is perpendicular to the radius, so the angular momentum is $mv\,r_{\min}$. Conservation of
angular momentum and of energy give
\[
mv_0b=mv\,r_{\min},\qquad
\tfrac12mv_0^2=\tfrac12mv^2+\frac{2Ze^2}{4\pi\varepsilon_0\,r_{\min}}.
\]
Divide the energy equation by $K=\tfrac12mv_0^2$, use $v=v_0b/r_{\min}$, and write $d$ for the head-on
distance of Eq.~\eqref{eq:nuc-closest}:
\[
1=\frac{b^2}{r_{\min}^2}+\frac{d}{r_{\min}}
\quad\Longrightarrow\quad
r_{\min}^2-d\,r_{\min}-b^2=0
\quad\Longrightarrow\quad
r_{\min}=\frac{d}{2}\left(1+\sqrt{1+\frac{4b^2}{d^2}}\,\right),
\]
taking the positive root. For $b=0$ this gives $r_{\min}=d$, the head-on result. For $b\gg d$ it gives
$r_{\min}\approx b$: a fast or distant particle hardly notices the nucleus. The full path is a hyperbola,
the repulsive counterpart of the conic orbits of \chref{ch:kepler}, and its deflection angle $\theta$
obeys $\tan(\theta/2)=d/(2b)$. That last result is stated without proof, because it needs the
equation of the orbit, which is beyond this book; Figure~\ref{fig:nuc-rutherford} was drawn from it.
\end{derivbox}

The deflection formula explains why back-scattering was so rare. A deflection of more than $90^\circ$
needs $\tan(\theta/2)>1$, that is, $b<d/2$. For \SI{7.7}{\MeV} $\alpha$ particles on gold, $d/2$ is
about \SI{15}{\fm}, so each gold nucleus presents a ``target'' of radius \SI{15}{\fm} inside an atom about
\SI{1.4e5}{\fm} in radius. Only a tiny fraction of the foil's area is covered by such targets.

\begin{exbox}{An off-centre approach}\label{ex:nuc-offcentre}
A \SI{7.7}{\MeV} $\alpha$ particle approaches a gold nucleus ($Z=79$) with impact parameter
$b=\SI{16}{\fm}$. Find its distance of closest approach and its deflection.

\emph{Solution.} The head-on distance is $d=2(79)(1.44)/7.7=\SI{29.5}{\fm}$. Then
\[
r_{\min}=\frac{29.5}{2}\left(1+\sqrt{1+\frac{4(16)^2}{29.5^2}}\,\right)=14.8\,(1+1.474)=\SI{36.6}{\fm}.
\]
The deflection is $\theta=2\tan^{-1}\!\big(29.5/(2\times16)\big)=2\tan^{-1}(0.923)=85^\circ$: this
particle leaves almost at right angles to its arrival, as the $b=\SI{16}{\fm}$ path in
Figure~\ref{fig:nuc-rutherford} shows, yet it never comes nearer than \SI{37}{\fm}.
\end{exbox}

\subsection{The nuclear radius}
To measure the actual size, the projectiles must get close enough to feel the nuclear force. Then the
scattering departs from the pure-Coulomb prediction, and the energy or angle at which that happens reveals
the radius. Experiments with $\alpha$ particles of higher energy, and especially with fast electrons
(which feel only the electric charge and so map out where the protons are), show that the nuclear
radius $R$ grows as the cube root of the mass number:
\begin{keybox}{Nuclear radius}
\begin{equation}\label{eq:nuc-radius}
R = R_0\,A^{1/3},\qquad R_0\approx\SI{1.2}{\fm}.
\end{equation}
This is an empirical rule, accurate to about ten percent for $A\gtrsim10$. $R$ is the radius of an
equivalent uniform sphere; real nuclei have no sharp edge, and their density falls from its central
value to near zero over about \SI{2}{\fm}.
\end{keybox}
The cube root is not an accident; it is the signature of constant density.

\begin{derivbox}{Why \texorpdfstring{$R\propto A^{1/3}$}{R proportional to A to the 1/3} means constant density}
If every nucleon occupies roughly the same volume $v$ and nucleons pack without overlapping, like marbles
in a bag, the nuclear volume is $V=Av$. A sphere has $V=\tfrac43\pi R^3$, so $R^3\propto A$ and
$R\propto A^{1/3}$. Conversely, given $R=R_0A^{1/3}$, the mass density of a nucleus of $A$ nucleons, each
of mass $m_{\text{nucleon}}$, is
\[
\rho=\frac{A\,m_{\text{nucleon}}}{\tfrac43\pi R_0^3A}=\frac{3\,m_{\text{nucleon}}}{4\pi R_0^3},
\]
in which $A$ has cancelled. With $m_{\text{nucleon}}\approx\SI{1}{\amu}=\SI{1.66e-27}{kg}$ and
$R_0=\SI{1.2e-15}{m}$,
\[
\rho\approx\frac{3(\num{1.66e-27})}{4\pi(\num{1.2e-15})^3}\approx\SI{2.3e17}{kg.m^{-3}}.
\]
\end{derivbox}

\begin{keybox}{Nuclear matter has one density}
Every nucleus, from carbon to uranium, has the same density, about \SI{2.3e17}{kg.m^{-3}}, which is
$\num{2e14}$ times that of water. A teaspoon (\SI{5}{mL}) of nuclear matter would have a mass of about a
billion tonnes. Nucleons behave like the molecules of an incompressible liquid drop, and this
\emph{liquid-drop} picture will explain the shape of the binding-energy curve in
Section~\ref{sec:nuc-curve}.
\end{keybox}

\begin{exbox}{Radii of helium-4 and uranium-238}\label{ex:nuc-radii}
Find the radii of $\iso{4}{He}$ and $\iso{238}{U}$, and the ratio of their volumes.

\emph{Solution.} $R_{\text{He}}=1.2\,(4)^{1/3}=\SI{1.90}{\fm}$ and
$R_{\text{U}}=1.2\,(238)^{1/3}=\SI{7.44}{\fm}$. The volume ratio is
$(7.44/1.90)^3=238/4=59.5$, the ratio of the mass numbers, as constant density requires. Uranium
contains 59.5 times as many nucleons but is only 3.9 times wider.
\end{exbox}

\begin{trapbox}{the radius does not scale with the mass number}
Doubling $A$ multiplies the radius by $2^{1/3}=1.26$, not by 2. It is the \emph{volume} that is
proportional to $A$. Ratios avoid $R_0$ entirely: $R_1/R_2=(A_1/A_2)^{1/3}$.
\end{trapbox}

\begin{remarkbox}{Neutron stars}
When the core of a massive star collapses, electrons are squeezed into protons, turning them into
neutrons (the reaction $p+e^-\to n+\nu_e$, where $\nu_e$ is the neutrino of \chref{ch:decay}), and the
core becomes in effect one gigantic nucleus, held together by gravity rather than by the nuclear
force. At nuclear density, 1.4 solar masses fits in a sphere only about \SI{14}{km} in radius (Practice
Problem~\ref{pp:nuc-neutronstar}). The estimate ignores the extra compression by gravity, and real
neutron stars are about \SI{12}{km} in radius. Some spin hundreds of times per second, and their
sweeping radio beams are observed as pulsars.
\end{remarkbox}

\section{Mass, Energy and the Atomic Mass Unit}\label{sec:nuc-massenergy}

The size of a nucleus came from mechanics. Its binding needs one idea that is new to this book. Einstein's
mass--energy relation states that mass is a form of energy: a body at rest whose internal energy
changes by $\Delta E$ changes its mass by
\[
\Delta m=\frac{\Delta E}{c^2},
\]
where $c=\SI{2.998e8}{m.s^{-1}}$ is the speed of light. Equivalently, a body of mass $m$ at rest has
\emph{rest energy} $mc^2$. This relation comes from special relativity, which is outside this book, so
it is used here as an experimentally established fact. Nuclear physics itself is among its most
precise tests: the mass differences measured with mass spectrometers agree with the energies released
in nuclear reactions, measured calorimetrically or from the speeds of the products.

Why did Volume~I never need it? For chemical reactions the change in mass is far too small to weigh, about
one part in $10^{10}$ (a few electron-volts released by molecules whose rest energy is tens of
\si{\giga\electronvolt}). For nuclear reactions it is about one part in a thousand, easily measured. This
is the first sign that nuclear energies are a million times larger than chemical ones.

Nuclear masses are quoted in \emph{unified atomic mass units} (\si{\amu}), defined so that one neutral
$\iso{12}{C}$ atom has mass exactly \SI{12}{\amu}. Thus $\SI{1}{\amu}=\SI{1.66054e-27}{kg}$, close to the
mass of one nucleon. Its rest energy is
\[
(\SI{1.66054e-27}{kg})(\SI{2.99792e8}{m.s^{-1}})^2=\SI{1.4924e-10}{J}=\SI{931.494}{\MeV}.
\]
\begin{keybox}{The conversion factor}
\[
(\SI{1}{\amu})\,c^2=\SI{931.494}{\MeV}.
\]
A mass difference measured in \si{\amu}, multiplied by $931.5$, gives the energy in \si{\MeV} directly.
Masses are often quoted in \si{\MeV}$/c^2$ (the electron's mass is $\SI{0.511}{\MeV}/c^2$), so that $c^2$
cancels silently in calculations.
\end{keybox}

\section{Mass Defect and Binding Energy}\label{sec:nuc-binding}

Now weigh a helium-4 atom and compare it with the total mass of the two hydrogen atoms and two
neutrons it could be built from. With the atomic masses $M(\iso{1}{H})=\SI{1.007825}{\amu}$ (the hydrogen
\emph{atom}, a proton plus an electron) and $m_n=\SI{1.008665}{\amu}$ (the neutron),
\[
2\,M(\iso{1}{H})+2\,m_n=2(1.007\,825)+2(1.008\,665)=\SI{4.032980}{\amu},\qquad
M(\iso{4}{He})=\SI{4.002603}{\amu}.
\]
The assembled atom is \emph{lighter} by \SI{0.030377}{\amu} (Figure~\ref{fig:nuc-massbalance}). That
missing mass left as energy, mostly as gamma rays (high-energy photons) and kinetic energy, when the
nucleons bound together; to pull them apart again, exactly that energy must be supplied.

\begin{figure}[htbp]
\centering
\tikzfig{nuc_massbalance}
\figcaption{fig:nuc-massbalance}{Weighing the parts against the whole. Two hydrogen atoms and two
neutrons (protons orange, neutrons grey) outweigh the helium-4 atom they can form, so the balance tips
towards the separated parts. The missing \SI{0.030377}{\amu} is the binding energy divided by $c^2$.}
\end{figure}

This is not peculiar to nuclei. \emph{Every} bound system weighs less than its separated parts, by its
binding energy divided by $c^2$; it is the mass--energy version of the negative total energy of a bound
orbit in \chref{ch:kepler}. A hydrogen atom is lighter than a free proton plus a free electron by
$\SI{13.6}{\electronvolt}/c^2$, about 1.4 parts in $10^{8}$ of its mass, far too little to weigh. Nuclear
binding is about a million times stronger per particle, so for nuclei the mass difference approaches one
percent and is measured routinely.

\begin{defbox}{Mass defect and binding energy}
For a nuclide $\nuc{A}{Z}{X}$ whose neutral atom has mass $M$, the \emph{mass defect} $\Delta m$ and the
\emph{binding energy} $\Eb$ are
\begin{equation}\label{eq:nuc-be}
\Delta m = Z\,M(\iso{1}{H})+(A-Z)\,m_n - M,\qquad \Eb=\Delta m\,c^2.
\end{equation}
$\Eb$ is the minimum energy needed to separate the nucleus completely into free protons and neutrons.
The \emph{binding energy per nucleon} $\Eb/A$ measures how tightly, on average, each nucleon is held.
\end{defbox}

For helium-4, $\Eb=(0.030\,377)(931.494)=\SI{28.30}{\MeV}$ and $\Eb/A=\SI{7.07}{\MeV}$.

\begin{trapbox}{electron bookkeeping}
Mass tables list \emph{atomic} masses, nucleus plus $Z$ electrons. Equation~\eqref{eq:nuc-be} uses the
hydrogen \emph{atom} mass $M(\iso{1}{H})$, not the bare proton mass, so $Z$ electron masses appear on both
sides and cancel automatically (the electrons' own binding energies, a few \si{\electronvolt} to
\si{\keV}, are negligible). Mixing a bare proton mass with an atomic mass leaves $Z$ electrons unaccounted
for. For uranium that is $92\times\SI{0.511}{\MeV}=\SI{47}{\MeV}$, a \num{2.6} percent error in $\Eb$.
\end{trapbox}

\begin{exbox}{Binding energy of oxygen-16}\label{ex:nuc-o16}
Given $M(\iso{16}{O})=\SI{15.994915}{\amu}$, find $\Eb$ and $\Eb/A$.

\emph{Solution.} $Z=8$ and $A-Z=8$, so
\[
\Delta m=8(1.007\,825)+8(1.008\,665)-15.994\,915=\SI{0.137005}{\amu}.
\]
Then $\Eb=0.137\,005\times931.494=\SI{127.6}{\MeV}$ and $\Eb/A=\SI{7.98}{\MeV}$, slightly above the
values for neighbouring nuclei. Section~\ref{sec:nuc-curve} explains why: oxygen-16 has filled shells
of both protons and neutrons.
\end{exbox}

\section{The Binding-Energy Curve}\label{sec:nuc-curve}

Binding energies computed in this way for many nuclides, divided by $A$ and plotted against $A$, give
Figure~\ref{fig:nuc-bea}. It is arguably the most important graph in nuclear physics, because the rest
of this volume reads its energetics from it.

\begin{figure}[htbp]
\centering
\tikzfig{nuc_bea}
\figcaption{fig:nuc-bea}{Binding energy per nucleon. The dots are computed with
Eq.~\eqref{eq:nuc-be} from measured atomic masses (2020 Atomic Mass Evaluation) for 35 nuclides; the
ringed ones are discussed in the text. The dashed curve is the liquid-drop formula of
Remark~\ref{rem:nuc-semf}. Notice the steep rise, the isolated high point of helium-4, the broad
maximum of about \SI{8.8}{\MeV} near iron-56 and nickel-62, and the slow decline beyond. Nuclei on either
side of the maximum release energy by moving towards it.}
\end{figure}

Its features, and their explanations, are these.
\begin{enumerate}[label=(\alph*)]
\item \textbf{A rough plateau near \SI{8}{\MeV} for $A\gtrsim20$.} If every nucleon attracted every other,
the total binding would grow as the number of pairs, which is about $A^2/2$, and $\Eb/A$ would grow in
proportion to $A$. It does not. The nuclear force has a short range, about \SIrange{2}{3}{\fm}, so each
nucleon binds only to its immediate neighbours. This \emph{saturation} makes $\Eb$ proportional to $A$,
just as the latent heat of a liquid is proportional to its mass.
\item \textbf{A steep rise for light nuclei.} A small nucleus is mostly surface, and a surface nucleon
has fewer neighbours to bind to. The surface area grows as $R^2\propto A^{2/3}$, so the fraction of
nucleons at the surface goes as $A^{2/3}/A=A^{-1/3}$ and shrinks with size.
\item \textbf{Spikes at $\iso{4}{He}$, $\iso{12}{C}$ and $\iso{16}{O}$.} Nucleons fill quantised energy
levels (shells), much as electrons do in atoms, and filled shells are extra stable. Helium-4, at
\SI{7.07}{\MeV} per nucleon against \SI{2.57}{\MeV} for helium-3, is so tightly bound that heavy nuclei
emit it as a ready-made unit: the $\alpha$ particle of \chref{ch:decay}.
\item \textbf{A maximum near $A\approx56$ to $62$}: $\Eb/A=\SI{8.79}{\MeV}$ for both $\iso{56}{Fe}$ and
$\iso{62}{Ni}$ (nickel-62 is higher by a few \si{\keV}).
\item \textbf{A slow decline for heavy nuclei}, to \SI{7.57}{\MeV} for $\iso{238}{U}$. Every proton repels
every other proton, not just its neighbours, so the electric energy grows as the number of proton pairs
divided by the size, $Z^2/R\propto Z^2/A^{1/3}$, faster than the saturated nuclear attraction.
\end{enumerate}

The curve converts at once into a rule for the energy of any nuclear process.
\begin{keybox}{Energy from rearranging nucleons}
In a nuclear reaction the numbers of protons and neutrons are conserved, so the energy released, called
the $Q$-value, equals the increase in total binding energy:
\begin{equation}\label{eq:nuc-qb}
Q=\sum E_{\mathrm{b},\text{products}}-\sum E_{\mathrm{b},\text{reactants}}.
\end{equation}
Free nucleons have $\Eb=0$. A positive $Q$ means energy is released, as kinetic energy of the products
or as gamma rays. Splitting a heavy nucleus (\emph{fission}) and joining two light ones
(\emph{fusion}) both move nucleons \emph{up} the curve towards the peak, so both release energy. Iron and
nickel are nuclear ``ash'': neither process can extract energy from them.
\end{keybox}
Equation~\eqref{eq:nuc-qb} follows from Eq.~\eqref{eq:nuc-be}: the $Z\,M(\iso{1}{H})$ and $(A-Z)\,m_n$
terms are the same on both sides because protons and neutrons are conserved, so the difference of
binding energies equals the difference of atomic masses times $c^2$, which is the energy released.
\chref{ch:decay} uses the mass form for decays, and \chref{ch:reactions} uses both forms for fission and
fusion. (In $\beta$ decay a neutron turns into a proton, so the numbers of protons and neutrons are
\emph{not} separately conserved; \chref{ch:decay} handles that case with atomic masses directly.)

\begin{exbox}{A fission estimate}\label{ex:nuc-fission}
A nucleus with $A=240$ and $\Eb/A=\SI{7.6}{\MeV}$ splits into two equal fragments with
$\Eb/A=\SI{8.5}{\MeV}$. Estimate the energy released.

\emph{Solution.} By Eq.~\eqref{eq:nuc-qb},
$Q=2(120)(8.5)-240(7.6)=2040-1824=\SI{216}{\MeV}$. A gain of only \SI{0.9}{\MeV} per nucleon, over 240
nucleons, gives about \SI{200}{\MeV} per fission. A typical chemical reaction releases a few
electron-volts per molecule, some fifty million times less.
\end{exbox}

\begin{trapbox}{per nucleon versus total}
$Q$ is a difference of \emph{total} binding energies. Subtracting the values of $\Eb/A$ directly gives
the gain per nucleon (\SI{0.9}{\MeV} above), not the energy released; multiply each $\Eb/A$ by its own
$A$ first.
\end{trapbox}

\begin{remarkbox}{The semi-empirical mass formula}\label{rem:nuc-semf}
In 1935 Carl Friedrich von Weizs\"acker turned the liquid-drop picture into a formula:
\[
\Eb(A,Z)=a_VA-a_SA^{2/3}-a_C\frac{Z(Z-1)}{A^{1/3}}-a_A\frac{(A-2Z)^2}{A}\pm\delta,
\]
with \emph{volume}, \emph{surface}, \emph{Coulomb}, \emph{asymmetry} and \emph{pairing} terms. The
constants are fitted to measured masses; one common set, used for the figures here, is $a_V=15.75$,
$a_S=17.8$, $a_C=0.711$ and $a_A=\SI{23.7}{\MeV}$. The first three terms were explained above: volume
binding from saturation, the surface deficit, and the repulsion between all $Z(Z-1)/2$ pairs of protons.
The asymmetry term reflects the \emph{Pauli exclusion principle} (no two identical nucleons may occupy
the same quantum state): neutrons and protons fill separate ladders of energy levels, so an excess of
either forces nucleons onto higher rungs. The small pairing term $\delta$ adds binding when $Z$ and $A-Z$
are both even and subtracts it when both are odd.

The dashed curve in Figure~\ref{fig:nuc-bea} is this formula, without the pairing term, with $Z$ at each
$A$ set to the most stable value. Apart from the lightest nuclei, where shell effects dominate, it
tracks the measured values to within two or three percent; it misses the helium-4 spike entirely,
because a five-parameter liquid drop knows nothing of shells. Figure~\ref{fig:nuc-semfterms} shows how
each term carves the curve. Setting $\partial\Eb/\partial Z=0$ at fixed $A$ predicts the most stable $Z$
for each $A$, the ``valley of stability'' of \chref{ch:decay}.
\end{remarkbox}

\begin{figure}[htbp]
\centering
\tikzfig{nuc_semfterms}
\figcaption{fig:nuc-semfterms}{Building the binding-energy curve term by term (liquid-drop formula,
per nucleon). Start from the constant volume term, which saturation predicts; subtract the surface
deficit, large for small nuclei; subtract the Coulomb repulsion, large for heavy nuclei; subtract the
neutron--proton asymmetry. What remains (blue) follows the measured values (dots) and peaks near
$A\approx60$, because the surface and Coulomb losses pull in opposite directions.}
\end{figure}

\section{The Nuclear Force and Separation Energies}\label{sec:nuc-force}

Binding requires a force that overpowers the electric repulsion of protons squeezed to within a few
femtometres of each other. Two protons \SI{2}{\fm} apart repel with a force of
$(\SI{1.44}{\MeV.\fm})/(\SI{2}{\fm})^2\approx\SI{0.36}{\MeV/\fm}$, which is about \SI{58}{N} acting on a
particle of mass \SI{1.7e-27}{kg}. The force that holds them anyway, the \emph{strong nuclear force},
has four features, each visible in the data already met.
\begin{itemize}
\item It is attractive, and near \SI{1}{\fm} its potential energy is tens of \si{\MeV} in magnitude,
far larger than the electric potential energy of two protons there (\SI{1.44}{\MeV}).
\item It has a short range, becoming negligible beyond about \SI{3}{\fm}. This is the saturation of
feature~(a).
\item It becomes strongly repulsive below about \SI{0.5}{\fm}, so nucleons cannot be squeezed together.
This is why nuclear density is constant.
\item It is \emph{charge-independent}: the proton--proton, neutron--neutron and neutron--proton nuclear
forces are essentially equal.
\end{itemize}
There is no simple formula for this force, unlike Coulomb's or Newton's law; it is a residue of the
deeper force between the quarks inside nucleons. Figure~\ref{fig:nuc-nnforce} sketches its potential
energy.

\begin{figure}[htbp]
\centering
\tikzfig{nuc_nnforce}
\figcaption{fig:nuc-nnforce}{Left: the potential energy of two nucleons against their separation
(schematic, not a fitted curve): a repulsive core inside about \SI{0.5}{\fm}, a deep attractive well near
\SI{1}{\fm}, and almost nothing beyond \SI{2}{\fm} to \SI{3}{\fm}. Right: zoomed out, the electric
repulsion between two protons is far weaker at contact but falls off only as $1/r$, so it alone
survives at large distances. In a big nucleus every proton repels \emph{all} the others but attracts
only its neighbours.}
\end{figure}

\textbf{Why do nuclei need neutrons?} A proton and a neutron do bind, though only just: the deuteron
$\iso{2}{H}$ has $\Eb=\SI{2.22}{\MeV}$, the smallest of any nucleus. Two protons do not bind, and
neither, perhaps surprisingly, do two neutrons, although no electric repulsion acts between them. The
nuclear attraction depends on how the nucleons' spins are aligned, and the Pauli exclusion principle
forbids two identical nucleons from taking the most strongly bound arrangement; the details are beyond
this book. (A common explanation, that two protons fail to bind only because of their electric
repulsion, cannot be right, since two neutrons fail too.) In larger nuclei, neutrons supply nuclear
attraction without adding any electric repulsion, so they act as the glue that holds many protons
together. The glue has a limit. Neutrons obey the exclusion principle, so adding too many forces them
onto ever higher energy levels, until converting a neutron into a proton becomes energetically
favourable. A stable nucleus is a compromise between these two effects, which is the subject of the
stability discussion in \chref{ch:decay}.

The average $\Eb/A$ hides individual variation. The energy needed to remove just one nucleon is the
\emph{separation energy}.
\begin{defbox}{Separation energies}
The \emph{neutron separation energy} $S_n$ of a nuclide with mass number $A$ and atomic number $Z$ is the
energy needed to remove one neutron:
\begin{equation}\label{eq:nuc-sn}
S_n=\big[M(A-1,Z)+m_n-M(A,Z)\big]c^2=\Eb(A,Z)-\Eb(A-1,Z),
\end{equation}
where $M(A,Z)$ is the atomic mass. The \emph{proton separation energy} is
$S_p=\big[M(A-1,Z-1)+M(\iso{1}{H})-M(A,Z)\big]c^2$; using the hydrogen-atom mass keeps the electrons
balanced, since the daughter atom has one electron fewer.
\end{defbox}
The second form of Eq.~\eqref{eq:nuc-sn} follows by writing both atomic masses with
Eq.~\eqref{eq:nuc-be}: everything cancels except the two binding energies. Separation energies are
large for nuclei with filled shells, just as ionisation energies are large for the noble gases.

\begin{exbox}{Removing one neutron from oxygen-16}\label{ex:nuc-sn}
Given $M(\iso{15}{O})=\SI{15.003066}{\amu}$ and $M(\iso{16}{O})=\SI{15.994915}{\amu}$, find $S_n$ for
$\iso{16}{O}$ and compare it with $\Eb/A$.

\emph{Solution.} $S_n=[15.003\,066+1.008\,665-15.994\,915](931.494)=\SI{15.66}{\MeV}$, about twice the
average $\Eb/A=\SI{7.98}{\MeV}$ of Example~\ref{ex:nuc-o16}. Removing a neutron from a filled shell is
expensive. Check with the second form of Eq.~\eqref{eq:nuc-sn}:
$\Eb(\iso{16}{O})-\Eb(\iso{15}{O})=127.62-111.96=\SI{15.66}{\MeV}$, the same answer.
\end{exbox}

\begin{trapbox}{one nucleon is not the average nucleon}
$\Eb/A$ is an average over all nucleons; the cost of removing the \emph{last} one, $S_n$ or $S_p$, can
be twice as large (filled shell) or much smaller (a lone nucleon outside a filled shell). A question that
asks for the energy to remove one neutron wants $S_n$, not $\Eb/A$.
\end{trapbox}

With the size, mass and binding of a nucleus in hand, the next question is what happens when a nucleus
is \emph{not} as tightly bound as it could be. \chref{ch:decay} shows that such a nucleus can lower its
energy by emitting a particle, and uses the atomic masses of this chapter to find how much energy each
decay releases.

\section{Summary}\label{sec:nuc-summary}
A nucleus is a dense drop of protons and neutrons whose mass is less than the mass of its parts; the
missing mass is the binding energy, and the curve of binding energy per nucleon decides which way
nuclear energy flows.

\begin{itemize}
\item Notation $\nuc{A}{Z}{X}$: $Z$ protons, $A$ nucleons, $A-Z$ neutrons. Isotopes share $Z$, isobars
share $A$, isotones share $A-Z$.
\item Closest approach of an $\alpha$ particle: head-on, $d=2Ze^2/(4\pi\varepsilon_0K)$ with
$e^2/4\pi\varepsilon_0=\SI{1.44}{\MeV.\fm}$; at impact parameter $b$, energy and angular momentum
conservation give $r_{\min}=\tfrac{d}{2}\big(1+\sqrt{1+4b^2/d^2}\big)$.
\item $R=R_0A^{1/3}$ with $R_0\approx\SI{1.2}{\fm}$, so nuclear density is the same for all nuclei,
about \SI{2.3e17}{kg.m^{-3}}.
\item Mass is energy: $(\SI{1}{\amu})c^2=\SI{931.494}{\MeV}$.
\item $\Delta m=Z\,M(\iso{1}{H})+(A-Z)\,m_n-M$ with atomic masses; $\Eb=\Delta m\,c^2$.
\item $\Eb/A$ rises steeply, peaks at about \SI{8.8}{\MeV} near iron-56 and nickel-62, and declines slowly
to \SI{7.6}{\MeV} for uranium; saturation, surface, Coulomb and asymmetry effects explain the shape.
\item $Q=\sum E_{\mathrm b,\text{products}}-\sum E_{\mathrm b,\text{reactants}}$: fission and fusion both
release energy by moving towards the peak.
\item The nuclear force is strong, short-ranged, repulsive at very short range and charge-independent.
\item $S_n=\Eb(A,Z)-\Eb(A-1,Z)$; separation energies are large for filled shells.
\end{itemize}

\section{Problem Types}\label{sec:nuc-types}
Five types cover what examiners ask about the nucleus itself. All of them reduce to two tools: the
cube-root rule for sizes, and mass differences times $931.5$ for energies.

\archq{Nuclear size and density}
{Radii, volumes or density; ``a nucleus whose radius is $n$ times that of \dots''.}
{Use $R=R_0A^{1/3}$. Ratios avoid $R_0$: $R_1/R_2=(A_1/A_2)^{1/3}$, and volume ratios equal
mass-number ratios. Density is the same for every nucleus.}
{Scaling the radius in proportion to $A$, or to $\sqrt{A}$, instead of $A^{1/3}$.}
{The radius of an $\iso{27}{Al}$ nucleus is \SI{3.6}{\fm}. The radius of a $\iso{125}{Te}$ nucleus is
approximately:}
{\SI{16.7}{\fm}}
{\SI{6.0}{\fm}}
{\SI{7.7}{\fm}}
{\SI{3.6}{\fm}\archans{B}{$R_{\text{Te}}=3.6\,(125/27)^{1/3}=3.6\times\tfrac53=\SI{6.0}{\fm}$.
Option (A) scales $R$ in proportion to $A$; (C) scales it as $\sqrt{A}$; (D) confuses equal density with
equal size.}}

\archq{Distance of closest approach}
{An $\alpha$ particle or proton of given energy approaches a nucleus; ``how close does it get'', with or
without an impact parameter.}
{Head-on: $K=q_1q_2/(4\pi\varepsilon_0 d)$, with $e^2/4\pi\varepsilon_0=\SI{1.44}{\MeV.\fm}$. Off-centre:
conserve energy \emph{and} angular momentum about the nucleus, $v_0b=v\,r_{\min}$, and solve
$r_{\min}^2-d\,r_{\min}-b^2=0$.}
{Forgetting that the $\alpha$ particle carries charge $2e$; or, off-centre, assuming the speed is zero
at closest approach (it is not: only the radial velocity vanishes).}
{A \SI{4.0}{\MeV} $\alpha$ particle approaches a gold nucleus ($Z=79$) with an impact parameter equal
to its head-on distance of closest approach $d$. Its distance of closest approach is about:}
{\SI{57}{\fm}}
{\SI{28}{\fm}}
{\SI{80}{\fm}}
{\SI{92}{\fm}\archans{D}{$d=2(79)(1.44)/4.0=\SI{56.9}{\fm}$, and with $b=d$,
$r_{\min}=\tfrac{d}{2}(1+\sqrt5)=1.618\,d=\SI{92}{\fm}$. Option (A) is the head-on value, which assumes
the speed vanishes at closest approach; (B) uses charge $e$ for the $\alpha$ particle; (C) is
$\sqrt2\,d$, a guess that adds $b$ and $d$ at right angles.}}

\archq{Mass defect and binding energy}
{Atomic masses are given; asked for $\Delta m$, $\Eb$, $\Eb/A$, ``which is more stable'', or the
minimum energy to break a nucleus into nucleons.}
{Count $Z$ and $A-Z$. Form $\Delta m=Z\,M(\iso{1}{H})+(A-Z)\,m_n-M$ with the hydrogen-\emph{atom} mass;
$\Eb=931.5\,\Delta m$ (\si{\MeV}, with $\Delta m$ in \si{\amu}); divide by $A$ to compare stability.}
{Using the bare proton mass with an atomic mass, which leaves $Z$ electrons unaccounted for.}
{Using $M(\iso{12}{C})=\SI{12.000000}{\amu}$, $M(\iso{1}{H})=\SI{1.007825}{\amu}$ and
$m_n=\SI{1.008665}{\amu}$, the binding energy per nucleon of carbon-12 is:}
{\SI{7.42}{\MeV}}
{\SI{15.4}{\MeV}}
{\SI{7.68}{\MeV}}
{\SI{92.2}{\MeV}\archans{C}{$\Delta m=6(1.007\,825)+6(1.008\,665)-12=\SI{0.098940}{\amu}$, so
$\Eb=\SI{92.16}{\MeV}$ and $\Eb/A=\SI{7.68}{\MeV}$. Option (A) uses the bare proton mass
(\SI{1.007276}{\amu}), losing six electrons; (B) divides by $Z$ instead of $A$; (D) is the total binding
energy, not yet divided by $A$.}}

\archq{Energy released, from the binding-energy curve}
{Values of $\Eb/A$ (or a graph) are given instead of masses; a nucleus splits or nuclei combine.}
{Convert each $\Eb/A$ to a total, $\Eb=A\times(\Eb/A)$, then $Q=\sum E_{\mathrm b,\text{products}}
-\sum E_{\mathrm b,\text{reactants}}$. Free nucleons have $\Eb=0$; positive $Q$ means energy released.}
{Subtracting the $\Eb/A$ values directly, which gives the gain per nucleon, not the energy released.}
{A nucleus with $A=180$ and $\Eb/A=\SI{8.0}{\MeV}$ splits into two equal fragments, each with
$\Eb/A=\SI{8.5}{\MeV}$. The energy released is:}
{\SI{90}{\MeV}}
{\SI{0.5}{\MeV}}
{\SI{45}{\MeV}}
{\SI{1530}{\MeV}\archans{A}{$Q=2(90)(8.5)-180(8.0)=1530-1440=\SI{90}{\MeV}$. Option (B) is the gain per
nucleon; (C) counts only one fragment's gain; (D) is the total binding energy of the products, not the
difference.}}

\archq{Separation energies}
{The energy to remove \emph{one} neutron or proton, rather than to take the whole nucleus apart.}
{Write the removal as a reaction and find its mass deficit, Eq.~\eqref{eq:nuc-sn}. For a neutron add
$m_n$; for a proton add $M(\iso{1}{H})$ and use the atomic mass of the daughter, so the electrons balance.
Compare with $\Eb/A$ to judge whether the nucleus is especially tightly bound.}
{Answering with $\Eb/A$, the average, instead of the cost of the one nucleon removed.}
{Given $M(\iso{12}{C})=\SI{12.000000}{\amu}$, $M(\iso{11}{C})=\SI{11.011434}{\amu}$,
$m_n=\SI{1.008665}{\amu}$ and $M(\iso{1}{H})=\SI{1.007825}{\amu}$, the energy needed to remove one neutron
from carbon-12 is:}
{\SI{7.68}{\MeV}}
{\SI{18.7}{\MeV}}
{\SI{17.9}{\MeV}}
{\SI{92.2}{\MeV}\archans{B}{$S_n=[11.011\,434+1.008\,665-12.000\,000](931.494)=\SI{18.7}{\MeV}$. Option
(A) is the average $\Eb/A$; (C) adds the hydrogen-atom mass instead of the neutron mass; (D) is the
energy to take the whole nucleus apart.}}

\section{Practice Problems}\label{sec:nuc-practice}

Use $R_0=\SI{1.2}{\fm}$, $e^2/4\pi\varepsilon_0=\SI{1.44}{\MeV.\fm}$, $(\SI{1}{\amu})c^2=\SI{931.494}{\MeV}$,
$M(\iso{1}{H})=\SI{1.007825}{\amu}$ and $m_n=\SI{1.008665}{\amu}$ throughout. Answers, with brief working,
follow the list.

\begin{enumerate}[itemsep=3pt]
\item\label{pp:nuc-ge} The radius of a germanium nucleus is measured to be twice the radius of a
$\iso{9}{Be}$ nucleus. How many nucleons are in the germanium nucleus?
\item\label{pp:nuc-ual} Find the ratio of the nuclear radius of $\iso{238}{U}$ to that of $\iso{27}{Al}$.
\item\label{pp:nuc-earth} If the Earth (\SI{5.97e24}{kg}) were compressed to nuclear density, what would
its radius be?
\item\label{pp:nuc-closest} In the Geiger--Marsden experiments of 1913 the $\alpha$ particles had
\SI{7.7}{\MeV}. (a) Find their distance of closest approach to a gold nucleus ($Z=79$). (b) What minimum
kinetic energy would an $\alpha$ particle need to touch the gold nucleus, taking both radii from
$R=R_0A^{1/3}$ ($A=197$ for gold, $A=4$ for the $\alpha$ particle) and neglecting recoil?
\item\label{pp:nuc-neutronstar} A neutron star has a mass of 1.4 solar masses
($M_\odot=\SI{1.989e30}{kg}$) at nuclear density. Estimate its radius.
\item\label{pp:nuc-fe} Given $M(\iso{56}{Fe})=\SI{55.934936}{\amu}$, find the binding energy and the
binding energy per nucleon of iron-56.
\item\label{pp:nuc-mirror} Tritium ($\iso{3}{H}$, \SI{3.016049}{\amu}) and helium-3 ($\iso{3}{He}$,
\SI{3.016029}{\amu}) are \emph{mirror nuclei}: each turns into the other if every proton is swapped for a
neutron and vice versa. (a) Compute both binding energies. (b) Account for the difference by estimating
the electric potential energy $e^2/(4\pi\varepsilon_0R)$ of two protons separated by about the nuclear
radius for $A=3$.
\item\label{pp:nuc-li7} Find $\Eb/A$ for $\iso{7}{Li}$ (\SI{7.016003}{\amu}) and locate it on
Figure~\ref{fig:nuc-bea}.
\item\label{pp:nuc-deut} The deuteron $\iso{2}{H}$ has atomic mass \SI{2.014102}{\amu}. What is the
minimum gamma-ray energy that can split it into a proton and a neutron? (Neglect the recoil.)
\item\label{pp:nuc-dd} Using $\Eb/A=\SI{1.11}{\MeV}$ for $\iso{2}{H}$ and \SI{7.07}{\MeV} for $\iso{4}{He}$,
find the energy released when two deuterons fuse into helium-4.
\item\label{pp:nuc-y} A nucleus Y ($A=200$, $\Eb/A=\SI{7.8}{\MeV}$) splits into fragments with $A=110$
($\Eb/A=\SI{8.5}{\MeV}$) and $A=90$ ($\Eb/A=\SI{8.6}{\MeV}$). Find $Q$.
\item\label{pp:nuc-3alpha} In red giant stars three $\alpha$ particles fuse into $\iso{12}{C}$ (the
triple-alpha process). Using $\Eb(\iso{4}{He})=\SI{28.30}{\MeV}$ and $\Eb(\iso{12}{C})=\SI{92.16}{\MeV}$,
find the energy released.
\item\label{pp:nuc-fefe} Explain, without calculation, why stars cannot generate energy by fusing two
$\iso{56}{Fe}$ nuclei, and why the core of a massive star therefore collapses once it has turned to iron.
\item\label{pp:nuc-sphe} Using $M(\iso{3}{H})=\SI{3.016049}{\amu}$, $M(\iso{3}{He})=\SI{3.016029}{\amu}$ and
$M(\iso{4}{He})=\SI{4.002603}{\amu}$, find (a) the proton separation energy and (b) the neutron separation
energy of $\iso{4}{He}$. Compare both with $\Eb/A$ for helium-4.
\item\label{pp:nuc-snsum} Show that the energies needed to remove all the nucleons of a nucleus one at a
time add up to its total binding energy $\Eb$, whatever the order of removal.
\end{enumerate}

\subsection*{Answers}
{\small
\begin{enumerate}[itemsep=2pt]
\item $R\propto A^{1/3}$, so $(A/9)^{1/3}=2$ and $A=72$ (germanium-72).
\item $(238/27)^{1/3}=2.07$.
\item $R=\big(3M/4\pi\rho\big)^{1/3}=\big(3(\num{5.97e24})/(4\pi\times\num{2.3e17})\big)^{1/3}\approx
\SI{180}{m}$: the whole planet would fit inside a sports stadium.
\item (a) $d=2(79)(1.44)/7.7=\SI{30}{\fm}$. (b) The surfaces touch when the centres are
$1.2\,(197^{1/3}+4^{1/3})=6.98+1.90=\SI{8.88}{\fm}$ apart, so $K=2(79)(1.44)/8.88=\SI{25.6}{\MeV}$,
more than three times what natural $\alpha$ sources provide. This is why Rutherford's $\alpha$ particles
never revealed the size of the nucleus.
\item $M=1.4\times\SI{1.989e30}{kg}=\SI{2.78e30}{kg}$, and $R=(3M/4\pi\rho)^{1/3}\approx\SI{14}{km}$.
\item $\Delta m=26(1.007\,825)+30(1.008\,665)-55.934\,936=\SI{0.52846}{\amu}$; $\Eb=\SI{492.3}{\MeV}$;
$\Eb/A=\SI{8.79}{\MeV}$.
\item (a) $\Eb(\iso{3}{H})=\SI{8.48}{\MeV}$ and $\Eb(\iso{3}{He})=\SI{7.72}{\MeV}$; the difference is
\SI{0.76}{\MeV}. (b) $R=1.2\,(3)^{1/3}=\SI{1.73}{\fm}$, so $e^2/(4\pi\varepsilon_0R)=1.44/1.73=
\SI{0.83}{\MeV}$. The two nuclei differ only by swapping neutrons and protons; with a charge-independent
nuclear force, the difference in binding should be the electric repulsion of helium-3's two protons, and
the estimate confirms it.
\item $\Delta m=\SI{0.04213}{\amu}$, $\Eb=\SI{39.2}{\MeV}$, $\Eb/A=\SI{5.61}{\MeV}$, on the steep
rising part of the curve between helium-4 and carbon-12.
\item $\Eb=[1.007\,825+1.008\,665-2.014\,102](931.494)=\SI{2.22}{\MeV}$.
\item $Q=4(7.07)-2\times2(1.11)=28.28-4.44=\SI{23.8}{\MeV}$.
\item $Q=110(8.5)+90(8.6)-200(7.8)=935+774-1560=\SI{149}{\MeV}$.
\item $Q=92.16-3(28.30)=\SI{7.26}{\MeV}$ (\SI{7.27}{\MeV} with unrounded masses).
\item Iron sits at the peak of the $\Eb/A$ curve. Any product heavier than iron has a lower $\Eb/A$, so the
total binding would \emph{decrease}: fusing iron absorbs energy instead of releasing it. With no energy
source left, the core can no longer support itself against gravity and collapses, producing a supernova
and often a neutron star.
\item (a) $S_p=[3.016\,049+1.007\,825-4.002\,603](931.494)=\SI{19.8}{\MeV}$.
\newline (b) $S_n=[3.016\,029+1.008\,665-4.002\,603](931.494)=\SI{20.6}{\MeV}$. Both are nearly three times
$\Eb/A=\SI{7.07}{\MeV}$: helium-4, with filled shells of both protons and neutrons, is exceptionally
tightly bound.
\item Removing nucleons one at a time, the separation energies are $\Eb(A)-\Eb(A-1)$, then
$\Eb(A-1)-\Eb(A-2)$, and so on, ending with the binding energy of the last pair minus zero, where each
$\Eb$ is that of whichever nuclide remains at that step. The sum telescopes to $\Eb(A)$ for any order,
because a free nucleon has zero binding energy.
\end{enumerate}
}
